% 第三章：结论与展望
% 总结全文工作、指出局限性、展望未来方向，一般 1～2 页。
\chapter{总结与展望}

\section{工作总结}
% 逐条对应绪论中的研究内容，概述做了什么、得到了什么结论
本文围绕示例问题开展了研究，主要工作总结如下：
\begin{itemize}
  \item 对问题进行了形式化建模，明确了优化目标；
  \item 设计并实现了示例系统，验证了方案的可行性；
  \item 实验结果表明，本文方法在主要指标上优于基线方法。
\end{itemize}

\section{不足与展望}
% 坦诚指出工作的局限性，并提出后续可改进的方向
受时间与条件限制，本文工作仍存在不足：样本规模有限、
对某些边界情况的讨论不够充分。未来可以从以下方向继续改进：
扩大实验规模、引入更多约束条件、探索方法的可解释性。
